\documentclass{beamer}

\title{Introduzione a Git}
\author{Nome Cognome}
\institute{Università degli Studi di XYZ}
\date{\today}

\begin{document}

\begin{frame}
\titlepage
\end{frame}

\begin{frame}
\frametitle{Cos'è Git?}
\begin{itemize}
\item Git è un sistema di controllo versione distribuito.
\item Permette di tenere traccia delle modifiche apportate ai file di un progetto nel tempo.
\item È stato creato da Linus Torvalds nel 2005 per gestire lo sviluppo del kernel Linux.
\end{itemize}
\end{frame}

\begin{frame}
\frametitle{Come funziona Git?}
\begin{itemize}
\item Git tiene traccia delle modifiche creando una serie di \textbf{commit}.
\item Ogni commit rappresenta uno stato del progetto in un determinato momento.
\item I commit sono organizzati in una \textbf{storia} (o \textbf{log}) dei cambiamenti.
\item Git utilizza un sistema di \textbf{branching} e \textbf{merging} per gestire le diverse linee di sviluppo.
\end{itemize}
\end{frame}

\begin{frame}
\frametitle{Comandi Git}
Ecco alcuni dei comandi Git più comuni:
\begin{itemize}
\item \texttt{git init}: inizializza un nuovo repository Git.
\item \texttt{git add}: aggiunge un file alla \textbf{staging area}.
\item \texttt{git commit}: crea un nuovo commit con le modifiche nella staging area.
\item \texttt{git push}: carica i commit sul server remoto.
\item \texttt{git pull}: scarica i commit dal server remoto e li integra nel repository locale.
\item \texttt{git clone}: clona un repository remoto sul computer locale.
\end{itemize}
\end{frame}

\begin{frame}
\frametitle{Conclusioni}
In questa presentazione abbiamo visto:
\begin{itemize}
\item Cos'è Git e come funziona.
\item Alcuni dei comandi Git più comuni.
\item Come Git può essere utile nella gestione dei progetti software.
\end{itemize}
Grazie per l'attenzione!
\end{frame}

\end{document}